%% ============================================================
%% 分片 32：原文第 11 节「$A$-Branes and ${\cal D}$-Modules」（原文 glshort.tex 起于第 9938 行）
%% 由 tools/harvmac2tex.py 从 arXiv 源码（hep-th/0604151v3）自动生成。
%% 只译英文散文；公式、编号、\label、\cite 一律不改。
%% ============================================================

\subsection{The Canonical Coisotropic $A$-Brane}

Our starting point is the observation made in \cite{ref80}\ that for a
general symplectic manifold $X$ there may exist $A$-branes which
are supported not on a Lagrangian submanifold but on a more
general coisotropic submanifold. A coisotropic submanifold $Y$ of
a symplectic manifold $X$ is a submanifold  defined locally by
first-class constraints, in other words it is defined locally by
the vanishing of functions $f_i$, $i=1,\dots,d$ that are
Poisson-commuting.  The maximal number of independent
Poisson-commuting functions on $X$ is one half the dimension of
$X$, so the dimension of $Y$ is at least half the dimension of
$X$.  If the dimension of $Y$ is precisely half the dimension of
$X$, then $Y$ is Lagrangian, and the $A$-branes supported on $Y$
are the usual Lagrangian $A$-branes.  Coisotropic $A$-branes with
support above the middle dimension also exist, in general, though
the conditions for their existence are rather special.

We will need only  the case that $Y=X$ and the $A$-brane has rank
one. Thus its Chan-Paton bundle is a $U(1)$ bundle ${\cal L}\to
X$. The Chan-Paton gauge field takes values in the Lie algebra of
$U(1)$, which is a one-dimensional real vector space. So this
gauge field is locally an ordinary real one-form $A$ with
covariant derivative $D=d+A$; its curvature is defined as $F=dA$.
In a unitary representation of $U(1)$, the element $1$ of the Lie
algebra $\frak{u}(1)$ acts by the $1\times 1$ anti-hermitian
matrix $-in$ for some integer $n$, and parallel transport is
$\exp(in\int A)$.  We interpret the Chan-Paton bundle over $X$ as a
line bundle associated to the representation $n=1$.  For a
sigma-model map $\Phi:\Sigma\to X$, the factor in the worldsheet
path integral associated with the Chan-Paton bundle is the
holonomy of $\Phi^*(A)$ around $\partial\Sigma$ in the
representation with $n=1$, or
$\exp(i\oint_{\partial\Sigma}\Phi^*(A))$.

In order for a brane with $Y=X$ and Chan-Paton line bundle ${\cal
L}$ to be an $A$-brane, its curvature $F$ must satisfy
% 式 (11.1)
\setcounter{equation}{0}
\begin{equation}\label{eq:coisobrane}
(\omega^{-1}F)^2=-1.
\end{equation}
 That is, the $(1,1)$ tensor
$\omega^{-1}F$ is an almost complex structure. The forms $F$ and
$\omega$ are closed, and together with \eqref{eq:coisobrane}\ this implies
that the almost complex structure $\hat N=\omega^{-1}F$ is
integrable \cite{ref80}. It is easy to see that both $\omega$ and $F$
are of type $(2,0)\oplus(0,2)$ with respect to $\hat N$.
Furthermore, the complex form $\omega-iF$ is of type $(2,0)$ and
so is a holomorphic symplectic form on $X$. Thus a coisotropic
$A$-brane of this kind may exist only if $\omega$ is the real part
of a holomorphic symplectic form.  See also \cite{ref67}\ for a
construction of branes of this type in generalized complex
geometry.

We are interested in the $A$-model on $\MH(G,C)$ with the
symplectic structure $\omega=(\Im\tau)\, \omega_K$. Then there is
an obvious solution to eqn. \eqref{eq:coisobrane} : % 式 (11.2)
\setcounter{equation}{1}
\begin{equation}\label{eq:yetro}
F=
(\Im\tau)\,\omega_J,\quad \hat N=\omega_K^{-1}\omega_J= I.
\end{equation}

(There is also a second solution with the sign of $F$ and $\hat N$
reversed.) The two-form $\omega_J$ is exact, as we explained in
section 第~\ref{sec:mhhyper}~节, so $F$ is the curvature two-form of a connection
on a trivial line bundle ${\cal N}\to\MH(G,C)$. The connection
form is % 式 (11.3)
\setcounter{equation}{2}
\begin{equation}\label{eq:exactcc}
\alpha={\Im\tau\over {2\pi}}\int_C |d^2z|\,
\Tr \left(\phi_z\delta A_\bz +\phi_\bz \delta A_z\right)
\end{equation}
 and
$F=d\alpha$. Although a more general two-form % 式 (11.4)
\setcounter{equation}{3}
\begin{equation}\label{eq:lombo}
(\Im\tau)\left(a\omega_J+b\omega_I\right), \quad a^2+b^2=1,
\end{equation}

solves eqn. \eqref{eq:coisobrane}, it is not exact for $b\neq 0$, and for
general values of $\tau$ does not have periods which are integral
multiples of $2\pi$. So in general, this form is not the curvature
two-form of any connection, and cannot be used to construct an
$A$-brane.

We will call the $A$-brane on $\MH$ with $F=(\Im\tau)\,\omega_J$
the canonical coisotropic brane, or c.c. brane, for short. We also
sometimes denote it as ${\cal B}_{c.c.}$. If $G$ is
simply-connected, then so is $\MH(G,C)$, and the c.c. brane is
unique.  Otherwise, it is unique up to twisting its Chan-Paton
bundle by a flat line bundle.

Our first observation about the c.c. brane is that it is an
$(A,B,A)$-brane, i.e. it is an $A$-brane in complex structures $I$
and $K$ and a $B$-brane in complex structure $J$. To check that it
is an $A$-brane in complex structure $I$, we only need the
identity % [WARN] 公式 \eqn\eqref{eq:hombo} 无编号或重复
\begin{equation*}
\tag*{(11.5)}
(\omega_I^{-1}\omega_J)^2=-1.
\end{equation*}
 This means
that the c.c. brane obeys the conditions of a coisotropic brane in
complex structure $I$, just as it does in complex structure $K$.
To check that it is a $B$-brane in complex structure $J$, it is
sufficient to note that the curvature of the gauge field on the
brane is a multiple of  $\omega_J$, which has type $(1,1)$ in
complex structure $J$.

These formulas really only define the c.c. brane on the smooth
part of $\MH$.  However, the c.c. brane can be defined by a
boundary condition in the underlying gauge theory; see section
第~\ref{sec:xgoto}~节\ for details.  This statement is the quantum field theory
equivalent of saying that the construction works on the whole
``stack'' of $G$-bundles over $C$.

Next, we would like to understand topological open strings with
both endpoints on the c.c. brane. They are represented by
BRST-invariant boundary operator insertions.  In general, as we
sketched in \dunamp\ of section 第~\ref{sec:comptwo}~节, for any brane ${\cal
B}$, the $({\cal B},{\cal B})$ strings form an associative, but
not necessarily commutative, algebra. In mathematical terms, this
is the endomorphism algebra of the brane ${\cal B}$ regarded as an
object of the category of branes on $\MH(G,C)$. We would like to
compute this endomorphism algebra in the case of the c.c. brane.

As a first step, we will argue that worldsheet instantons do not
contribute to the computation. The boundary condition defining
coisotropic branes in the sigma-model \cite{ref80}\ says that strings
ending on a space-filling coisotropic brane such as we consider
here must obey % 式 (11.6)
\setcounter{equation}{5}
\begin{equation}\label{eq:melbo}
G_{IJ}\partial_1\Phi^J=-i
F_{IJ}\partial_0\Phi^J,
\end{equation}
 where $\partial_1$ and $\partial_0$ are
normal and tangential derivatives at the worldsheet boundary.
Imposing also the instanton equation (that is, the map $\Phi$
should be holomorphic in complex structure $K$), to relate
$\partial_1\Phi$ to $\partial_0 \Phi$, we get % 式 (11.7)
\setcounter{equation}{6}
\begin{equation}\label{eq:zimbolo}
(\omega-iF)\partial_0\Phi=0.
\end{equation}
 Since both $\omega$ and $F$ are
nondegenerate, the instanton must be constant on the boundary. But
then, by analyticity, it must be constant everywhere, hence
trivial.

Since there are no instanton corrections, anything interesting
will have to come from worldsheet perturbation theory.  But
worldsheet perturbation theory, to any order in $\hbar$, generates
corrections to sigma-model operators and couplings that are local
in the target space. Hence, at least to all orders in world-sheet
perturbation theory, it is very natural to consider boundary
observables that may be defined only locally in target space. Such
observables form a sheaf over the target space, so instead of an
algebra of observables, we can get, to all orders in perturbation
theory, something more powerful: a sheaf of algebras. (In the
closed-string case, such localization in target space has been
discussed in \cite{ref133,ref134,ref135}, interpreting earlier
mathematical work \cite{ref136}.) Along with analyzing this
sheaf in perturbation theory, we will discuss what happens in the
exact theory.

The  space of BRST-invariant boundary operators for an arbitrary
coisotropic $A$-brane $Y$ has been determined in the classical
limit in \cite{ref80,ref132}. For our purposes, it will be
sufficient to consider boundary operators of ghost number zero.
Then, in the case $Y=X$, such observables are simply holomorphic
functions on $X$, where $X$ is regarded as a complex manifold with
complex structure $\hat N$. Specializing to this case and
localizing in target space, we conclude that additively (ignoring
the ring structure) the sheaf of topological boundary operators
for the c.c. brane is the sheaf of functions on $\MH(G,C)$ which
are holomorphic in complex structure $I$. Note that while we are
working in the $A$-model in complex structure $K$, the complex
structure which is relevant here is $\hat N=I$. Locally the sheaf
of boundary operators has no higher cohomology and cannot be
deformed, so it cannot have quantum corrections. But globally it
has quantum corrections, which we will analyze.

The sheaf of holomorphic functions on $\MH$ has an obvious
commutative algebra structure, which is the correct one in the
classical limit. In the context of ``physical'' string theory, in
the presence of the gauge field strength $F$, the algebra of open
strings becomes noncommutative \cite{ref137}.

Let us first discuss informally how
this would work in the present context.  The usual noncommutative parameter, called $\theta$ in
\cite{ref137}, is the antisymmetric
part of $(g+F)^{-1}$.  In the present context, this becomes a multiple of $\omega_J^{-1}$.  $\omega_J$
is the real part of $\Omega_I$, the holomorphic $(2,0)$-form in complex structure $I$.  When we multiply
observables derived from functions
that are holomorphic in complex structure $I$ (which are the topological
observables), $\omega_J^{-1}$ can be replaced by $\Omega_I^{-1}$.  So the multiplication of such
observables  is deformed via the Poisson brackets derived from
$\Omega_I$, as we will find shortly.

In physical string theory, for any $F\not=0$, the multiplication of open strings contains noncommutative
phases coming from $(g+F)^{-1}$, but generally there are many other stringy effects.  A precise
description in terms of a noncommutative deformation of the algebra of functions on the target space
usually arises \cite{ref137}\ only when $|F|>>|g|$.
We are not in that limit in the present problem; rather, we have $|F|=|g|$.  However,
once we restrict to topological observables, we will get a simplification
similar to what usually occurs in physical string theory for $|F|>>|g|$.
In the context of the $B$-type topological
strings, this was first discussed in \cite{ref71}.  In essence, the elimination of stringy excitations that
in physical string theory occurs only if $|F|>>|g|$ arises here because of restricting to topological
observables.

Now we turn to a more precise analysis for the topological theory.
The action for the
$A$-model on a disk can be written in the form % 式 (11.8)
\setcounter{equation}{7}
\begin{equation}\label{eq:mortop}
S=\int_\Sigma \Phi^*(\omega-iF)+\int_\Sigma d^2\sigma \{Q,V\},
\end{equation}

where $Q$ is the BRST operator or topological supercharge. In our
case, the form $\omega-iF=-i \Omega_I \Im\tau$ is exact (and of
course this is always true locally). If we introduce the potential
$\varpi$ for $\Omega_I$, so that $\Omega_I=d\varpi$, then the
action takes the form % 式 (11.9)
\setcounter{equation}{8}
\begin{equation}\label{eq:ortop}
S=-i\,\Im\tau\int_{\partial\Sigma}\Phi^*(\varpi) +\int_\Sigma
d^2\sigma \{Q,V\}.
\end{equation}
 Further,  $\varpi$ is proportional to the
canonical holomorphic one-form on $\M_H(G,C)$: % 式 (11.10)
\setcounter{equation}{9}
\begin{equation}\label{eq:honop}
\varpi={1\over \pi}\int |d^2 z|\, \Tr \left(\phi_z\delta
A_\bz\right).
\end{equation}
  If we restrict to $T^*\M(G,C)$, we can be more
explicit.  If $q^\alpha$ are local holomorphic coordinates on
$\M(G,C)$ and $p_\alpha$ are linear functions on the fibers of
$T^*\M(G,C)$ that are canonically conjugate to the $\phi_\alpha$,
then $\varpi=\sum_\alpha \,p_\alpha\,dq^\alpha$ and $S=-i\,{\rm
Im}\,\tau\int p_\alpha\,dq^\alpha+\{Q,\dots\}$.  Thus, up to
$Q$-exact terms, the $A$-model action is the holomorphic version
of the quantum mechanical action for a particle on $\cM(G,C)$
(with zero Hamiltonian).

We undoubtedly lose some important structure by specializing to
$T^*\M\subset \MH$, since the c.c. brane is defined on all of
$\MH$, and even in the four-dimensional gauge theory, not just on
$T^*\M$.  But let us see what structure remains after we make this
specialization.

The form of the action in \eqref{eq:ortop}\ is standard, and the result is
familiar: to first order in ${1/ {\rm Im}\,\tau}$, the effect of
the perturbation is to deform the commutative algebra of
locally-defined holomorphic functions by the Poisson brackets
derived from this action.  So in perturbation theory in ${1/{\rm
Im}\,\tau}$, the c.c. strings (or more fastidiously, the $({\cal
B}_{c.c.},{\cal B}_{c.c.})$ strings) form a sheaf of
noncommutative algebras. But precisely what sheaf do we get?

A powerful tool is the $\CC^*$ group ${\cal U}$ that was
introduced at the end of section 第~\ref{sec:complexstr}~节\ and used at various
points in this paper. It acts by % 式 (11.11)
\setcounter{equation}{10}
\begin{equation}\label{eq:utyp}
\phi_z\to \lambda
\phi_z
\end{equation}
 and the associated grading on the space of observables was
used in section 第~\ref{sec:sdualhitfib}~节\ to analyze the action of $S$-duality
on zerobranes. $\varpi$ is of degree 1 under this scaling;
therefore,  ${\cal U}$ can be promoted to a symmetry of the
quantum theory if we assign degree $1$ to ${1/{\rm Im}\,\tau}$,
that is to Planck's constant $\hbar$ in the analog
quantum-mechanical problem.

At the classical level, the multiplication of observables
preserves the grading by ${\cal U}$; the product of observables of
degrees $m$ and $n$ is an observable of degree $m+n$. At the
quantum level, there can also be terms of lower degree in the
product, if they are accompanied by positive powers of  Planck's
constant. Thus the algebra of quantum observables is only
filtered, not graded.

In order to effectively use the $\Bbb{C}^*$ symmetry to analyze
the algebra, we need to consider observables that are defined in a
$\Bbb{C}^*$-invariant open set.  For this, we pick an open set
${{\cal V}}\subset {\EUM}$ and consider observables that are
defined on $T^*{{\cal V}}\subset \MH$.  Thus, we are only going to
partly sheafify the space of strings ending on the c.c. brane,
allowing only open sets in $\MH$ of this form. The advantage of
this partial sheafification is that it works nonperturbatively, as
we will see shortly.  A complete sheafification in which one
defines an associative algebra of observables for an arbitrary
open set ${\cal W}\subset \MH$ works fine to all orders in $1/{\rm
Im}\,\tau$, but it is not realistic to expect it to work in the
exact theory. For a general discussion of such matters, see
\cite{ref138}.

To understand why in general there is an obstruction to
sheafification of deformation quantization beyond perturbation
theory, consider the simple example of deformation quantization of
$\Bbb{C}^{2}$ with coordinates $x,y$ and symplectic form $dx\wedge
dy$. Deformation quantization in this case can be carried out
exactly using the Moyal-Wigner product\footnote{For reprints of their
papers and many other original papers, along with a survey, see
\cite{ref139}.} if one restricts oneself to
polynomial functions on $\Bbb{C}^{2}$. To sheafify this situation,
one picks an open set ${\cal W}\subset\Bbb{C}^{2}$, and tries to
define a Moyal-Wigner product of functions holomorphic in ${\cal
W}$. There is no problem to any finite order in $\hbar$, but the
fact that the Moyal-Wigner product involves an infinite sum causes
a problem in getting an exact result.  To see the problem, try to
evaluate at the origin the Moyal-Wigner product of the functions
$1/(x-a)$ and $1/(y-b)$, each of which are holomorphic in suitable
open sets in $\Bbb{C}^2$.

So we limit ourselves to partial sheafification in which we
consider an open set in $\MH$ of the form $T^*{{\cal V}}$, ${{\cal
V}}\subset\M$.  Moreover, we restrict ourselves to functions on
$T^*{{\cal V}}$ that have polynomial growth along the fibers of
the projection $T^*{{\cal V}}\to{{\cal V}}$.  The algebra of such
functions is generated by elements of degree less than or equal to
$1$, that is, by functions that are at most linear along the
fibers.

Elements of degree $0$ are simply holomorphic functions on ${{\cal
V}}$.  There is no possibility to deform the algebra of
multiplication of such functions, because there are no observables
in $T^*{{\cal V}}$ of negative degree.  As $\hbar$ has degree 1, a
correction to the multiplication of functions of degree 0 would
have to involve an observable of negative degree.

Elements of degree $1$ are linear functions on the fibers of the
cotangent bundle.  So if $q^\alpha$ are local coordinates on
${{\cal V}}$ and $p_\alpha$ are the conjugate linear functions on
the fibers, then an element of degree 1 takes the form
$\kappa=\sum_\alpha p_\alpha\, V^\alpha(q)$ and is in natural
correspondence with a holomorphic vector field $\sum_\alpha
V^\alpha{\partial\over\partial q^\alpha}$ on $\cM$. In multiplying
an object $f$ of degree $0$ and an object $\kappa$ of degree 1,
there is only ``room'' for a first order deformation -- since
terms of order $\hbar^2$ in the product would again multiply
operators of negative degree. The first order deformation is
characterized by the Poisson bracket % 式 (11.12)
\setcounter{equation}{11}
\begin{equation}\label{eq:hembo}
[\kappa,
f]=-i(\Im\tau)^{-1}\kappa(f).
\end{equation}

To proceed, we write $\hat w$ for the open string state
corresponding to a holomorphic function $w$ on $T^*{{\cal V}}$.
The formula \eqref{eq:hembo}\ determines the $\star$ product (that is, the
algebra of open string multiplication) of the ring elements $\hat
f$, $\hat\kappa$ corresponding to $f$ and $\kappa$, but not quite
uniquely, since (while preserving the way the algebra is filtered)
we could add to $\hat\kappa$ an observable of degree zero. If for
$\kappa=V^\alpha(q) p_\alpha$, we define $\hat\kappa =\hat
V^\alpha\star \hat p_\alpha$, then the algebra is uniquely
determined and is % 式 (11.13)
\setcounter{equation}{12}
\begin{equation}\label{eq:yembo}
\begin{aligned}\hat f\star\hat\kappa&
=\hat{f\kappa}\\ \hat\kappa\star \hat f&
 =\hat{f\kappa}-i(\Im\tau)^{-1}\hat{\kappa(f)}.\\\end{aligned}
\end{equation}

 Of course, $\hat{ f\kappa}$ and $\hat{\kappa(f)}$ are the ring
 elements that correspond to the functions
 $f\kappa$ and $\kappa(f)=V^\alpha\,\partial f/\partial
q^\alpha$. We still have not completely used the freedom to
redefine $\hat\kappa$; the remaining freedom is that for the
linear functions $p_\alpha$ on the cotangent bundle, we could
replace % 式 (11.14)
\setcounter{equation}{13}
\begin{equation}\label{eq:remob}
\hat p_\alpha\to \hat p_\alpha+\hat A_\alpha
\end{equation}

for some functions $A_\alpha(q)$.

When we multiply two functions $\kappa$ and $\kappa'$ that are
each of degree 1, something new happens. There is a first order
deformation, which is given again by the Poisson bracket, but
there now is also room for a second order deformation.  The
Poisson bracket of $p_\alpha $ and $p_\beta$ vanishes, but the
commutator $[\hat p_\alpha,\hat p_\beta]$ in the $\star$ algebra
might be $\hbar^2 F_{\alpha\beta}(q)$ for some functions
$F_{\alpha\beta}$. One can show using associativity of the algebra
that the two-form $F_{\alpha\beta}dq^\alpha\,dq^\beta$ is closed,
and hence can be set to zero (for a small open set ${{\cal V}}$)
by the transformation \eqref{eq:remob}.

The algebra associated with $T^*{{\cal V}}$ for a small open set
${\cal V}\subset \MH$ is therefore isomorphic simply to the
algebra of holomorphic differential operators on ${{\cal V}}$,
with functions on $\M$ acting by multiplication and
% 式 (11.15)
\setcounter{equation}{14}
\begin{equation}\label{eq:mebolo}
\hat p_\alpha=-{i\over {\rm
Im}\,\tau}{\partial\over\partial q^\alpha}.
\end{equation}

Now suppose that we cover $\M$ with small open sets ${{\cal
V}}^{(i)}$.  In overlap regions ${{\cal V}}^{(i)}\cap {{\cal
V}}^{(j)}$, the degree zero ring elements $\hat f$ must agree, but
the degree 1 elements $\hat p_\alpha^{(i)}$ and $\hat
p_\alpha^{(j)}$ may differ by a transformation \eqref{eq:remob}.  Since we
require $[\hat p_\alpha,\hat p_\beta]=0$ in each open set, they
must be related by % 式 (11.16)
\setcounter{equation}{15}
\begin{equation}\label{eq:hecox}
\hat p_\alpha^{(i)}=\hat
p_\alpha^{(j)}+{1\over {\rm Im}\,\tau}\hat{\partial s^{(ij)}\over
\partial q^\alpha}
\end{equation}
 for some local holomorphic functions $s^{(ij)}$ defined on
${{\cal V}}^{(i)}\cap {{\cal V}}^{(j)}$.   Consistency of this
relation implies that in triple overlaps the quantities
$c^{(ijk)}=s^{(ij)}+s^{(jk)}+s^{(ki)}$ are complex constants.
Moreover, they automatically obey a cocycle condition and define
an element of $H^2({\EUM},\Bbb{C})$.

If those constants take values in $2\pi\Bbb{Z}$, then the
quantities $\exp(-is^{(ij)})$ are the transition functions of a
complex line bundle ${\cal L}\to \MH$.  In this case, the sheaf of
algebras of open strings is the sheaf of differential operators
acting on sections of the line bundle ${\cal L}$.  We call this
sheaf of algebras ${\cal D}_{\cal L}$.

${\cal L}$ is not quite uniquely determined.  In general, the
sheaf of algebras ${\cal D}_{\cal L}$ is invariant under twisting
${\cal L}$ by a flat line bundle, since a flat line bundle has
constant transition functions that commute with the momenta $\hat
p_\alpha$.

There is no reason that the constants $c^{(ijk)}$ must be properly
quantized, that is, must take values in $2\pi\Bbb{Z}$.  For
example, beginning with a properly quantized case leading to the
sheaf ${\cal D}_{\cal L}$ of differential operators on a line
bundle ${\cal L}$, we can define a new sheaf ${\cal D}_{\cal
L^\gamma}$ by multiplying the $s^{(ij)}$ by a complex number
$\gamma$. If ${\cal L}$ is topologically nontrivial, a line bundle
${\cal L}^\gamma$ does not exist (except for integer or perhaps
rational values of $\gamma$) but the corresponding sheaf of
algebras ${\cal D}_{{\cal L}^\gamma}$ does exist.

Since $H^2({\EUM}(G,C),\Bbb{C})$ is one-dimensional, in this
specific example the most general sheaf of algebras that could
arise from this sort of construction is of the form ${\cal
D}_{{{\frak L}}^\gamma}$ where $\gamma$ is a complex number and
${{\frak L}}$ is the fundamental line bundle over $\EUM$, which we
loosely call the determinant line bundle. If ${\EUM}$ is replaced
by a more general space $X$, a slight generalization is possible
involving a product of complex powers of line bundles.

At any rate, as we will see, the basic case of the geometric
Langlands program involves the differential operators on a certain
ordinary line bundle.  But a generalization that we consider in
section 第~\ref{sec:genotics}~节\ does involve the complex powers of a line
bundle.

Notice that in the above analysis, there is no question of whether
perturbation theory converges, since everything is determined by
what happens in second order.  Had we tried to completely
sheafify, associating an algebra with every open set in $\MH$, we
would have had trouble with convergence of perturbation theory.

\par\medskip{\it Time Reversal Symmetry}

These somewhat abstract considerations do not quite determine the
sheaf of algebras of c.c. strings. It is the sheaf of differential
operators on some line bundle (or complex power of a line bundle),
but which one? To determine it, we will use time-reversal
symmetry.  It is natural to get an extra constraint from
time-reversal symmetry, because as we see in eqn. \eqref{eq:zongor}, this is
a symmetry of the twisted topological field theory precisely at
$t=1$, the value at which we are working in the present
discussion.

In gauge theory, time-reversal, which we  call ${\EUT}$, reverses
the time coordinate of the four-manifold $M$ while acting
trivially on the space coordinates.  If $M=\Bbb{R}\times I\times
C$, then ${\EUT}$ just acts as $-1$ on $\Bbb{R}$. In the effective
two-dimensional theory, ${\EUT}$ is an orientation-reversing
symmetry of the string worldsheet $\Sigma$, which maps the
boundary of $\Sigma$ to itself while reversing the orientation of
the boundary.  Since our considerations are local, we can assume
that $\Sigma$ or $M$ has such a symmetry.

Under ${\EUT}$, the bosons $(A,\phi)$ transform to
$({\EUT}^*A,-{\EUT}^*\phi)$, as we saw in eqn. \eqref{eq:yoro}.  The terms
in \eqref{eq:mortop}\ and \eqref{eq:ortop}\ that involve $\omega$, $F$, and $\varpi$
would be odd under reversal of worldsheet orientation if ${\EUT}$
were to map $\phi$ to ${\EUT}^*\phi$. But with the extra minus
sign, these terms are invariant.

Time-reversal reverses the order of operator insertions on the
worldsheet boundary, so it maps the algebra of c.c. strings to its
``opposite'' algebra, in which the elements are the same but
multiplication goes in the opposite order. Thus, to every element
$x$ of an algebra ${\eurm C}$, the opposite algebra ${\eurm
C}^{op}$ has an element $x^{op}$, with $(x\star
y)^{op}=y^{op}\star^{op} x^{op}$.  Here $\star $ is the
multiplication in ${\eurm C}$ and $\star^{op}$ is the
multiplication in ${\eurm C}^{op}$.

  
The opposite of the algebra of differential operators on a line
bundle $\cL$ over a complex  manifold $X$ is \cite{ref140}\ isomorphic
to the algebra of differential operators on $\cL^{-1}\otimes K_X$,
where $K_X$ is the canonical line bundle of $X$ (that is, the
bundle of holomorphic forms of top degree). To see this, let $s$
and $u$ be compactly supported (not holomorphic) sections,
respectively, of $\cL$ and $\cL^{-1}\otimes K_X$, and let ${\eurm
F}$ be a holomorphic differential operator acting on $\cL$.  Then
$s\, {{\eurm F}}u$ is a section of $K_X$. A section of $K_X$ is a
differential form of middle dimension, and can be integrated over
any real slice $Z$ of $X$. Moreover, by integrating by parts, we
can form a ``transpose'' operator ${{\eurm F}}^t$, acting on
sections of $\cL^{-1}\otimes K_X$, and obeying $\int_Z u\,{{\eurm
F}} s=\int_Z ({\eurm F}^t u)s$ for all $u$ and $s$. ${{\eurm
F}}^t$ is a differential operator with holomorphic coefficients,
just like ${{\eurm F}}$, and does not depend on the choice of $Z$.
The map ${{\eurm F}}\to {{\eurm F}}^t$ reverses the order of
multiplication, showing that the opposite sheaf of algebras to
${\cal D}_{\cal L}$ is ${\cal D}_{{\cal L}^{-1}\otimes K_X}$.

Time-reversal symmetry of the c.c. brane means that the sheaf of
algebras obtained from open strings ending on the c.c. brane must
be isomorphic to its opposite algebra.  This, plus the fact that
it is the sheaf of differential operators acting on some line
bundle, implies that this algebra must be ${\cal
D}_{K_{\EUM}^{1/2}}$, the sheaf of differential operators acting
on $K_{\EUM}^{1/2}$, the square root of the canonical bundle of
$\M$. $\M$ actually is a spin manifold, so $K_\M^{1/2}$ is an
ordinary line bundle.  The spin structure of $\M$ is not
necessarily unique, but since the sheaf of differential operators
acting on a line bundle is invariant under twisting by a flat line
bundle, the sheaf of algebras ${\cal D}_{K_\EUM^{1/2}}$ does not
depend on the choice of $K_{\EUM}^{1/2}$.

\par\medskip {\it Global Algebra of Topological Strings}

Now we want to discuss the {\it global} c.c. strings.

How do we get the algebra of global observables from the sheaf of
observables?  We simply take the global sections of the sheaf of
observables, or more generally, the cohomology of $\M$ with values
in this sheaf. To keep things simple, we will just consider global
sections, that is observables of ghost number or cohomological
degree zero (but any scaling degree under the $\Bbb{C}^*$ symmetry
${\cal U}$).

Classically, the sheaf of observables is just the sheaf of
holomorphic functions on $\MH$, and its global sections are simply
the global holomorphic functions on $\MH$.  They are the commuting
Hamiltonians of the integrable system, which we described in
section 第~\ref{sec:hitchfib}~节.

Likewise quantum mechanically, to get the global observables, we
must take the global sections of the sheaf of local observables.
But there is a quantum correction to the gluing law that defines
this sheaf, relative to what we have classically.  The quantum
correction entered our analysis explicitly in eqn. \eqref{eq:hecox}. So
there is potentially a quantum correction to the description of
the global observables.

Having understood that the sheaf of c.c. strings is the sheaf of
differential operators on $\MH$, we can understand concretely what
this correction means.  Suppose that we are given a holomorphic
function on $\MH$, for example a function of degree three that
once we pick local complex coordinates on $\M$ looks like
$w=\sum_{\alpha,\beta,\gamma} f^{\alpha\beta\gamma}(q)p_\alpha
p_\beta p_\gamma$.  To promote it to a global $({\cal
B}_{c.c.},{\cal B}_{c.c.})$ string, we need to find a third order
differential operator $\hat w$ on $\M$ whose ``symbol'' is $w$. In
other words, we want $\hat w=(-i/{\rm
Im}\,\tau)^3\,f^{\alpha\beta\gamma}\partial^3 /\partial
q^\alpha\partial q^\beta\partial q^\gamma+\dots$, where the
ellipses refer to a differential operator of order two or less.
Because of the absence of a canonical system of local complex
coordinates $q^\alpha$ on $\M$, such an operator is not unique
locally, and may not exist globally.

If a holomorphic differential operator $\hat w$ with symbol $w$
does exist, we say that the function $w$ can be ``quantized.'' The
commuting Hamiltonians on $\MH$ can all be quantized.  This was
first proved for functions of degree two by Hitchin \cite{ref141}\
and for functions of any degree by Beilinson and Drinfeld \cite{ref5};
this result is part of the rationale for the title of their
paper.\footnote{Their argument is based on two-dimensional conformal
field theory, more precisely on current algebra of the group $G$
at level $k=-h$, with $h$ the dual Coxeter number of $G$.} We will
give an argument that involves further use of the time-reversal
symmetry $\EUT$. This argument also shows that the resulting
differential operators commute, as first proved in the above-cited
references.

We use the fact that $({\cal B}_{c.c},{\cal B}_{c.c.})$ strings
correspond to boundary operators that can be inserted at a
boundary labeled by the c.c. brane.  Moreover, these boundary
observables can be derived from operators of the four-dimensional
gauge theory, which is a powerful tool because it has properties
of locality that are lost upon reduction to two dimensions.  This
will enable us to give a simple argument.  In fact, as we recall
from section 第~\ref{sec:hitchfib}~节, the commuting Hamiltonians take the form
$\int_C \alpha\,{\cal P}(\varphi)$, with ${\cal P}$ a homogeneous
gauge-invariant polynomial of some degree $k$, and $\alpha\in
H^1(C,K_C^{1-k})$.  To show that such an operator is
$Q$-invariant, it suffices to show that the local operator ${\cal
P}(\varphi)$ is $Q$-invariant.  We do this in section 第~\ref{sec:xgoto}~节,
simply by classifying the possible operators and using
time-reversal symmetry.  From this point of view, it is clear that
the differential operators obtained by quantizing two classical
expressions $\int_C \alpha \,{\cal P}(\varphi)$ and $\int_C\tilde
\alpha \,\tilde {\cal P}(\varphi)$ commute, since one can take
$\alpha$ and $\tilde\alpha$ to have disjoint support.

To illustrate the nontrivial nature of the fact that the classical
holomorphic functions on $\M$ can be quantized to get differential
operators on $K_\M^{1/2}$, let us note \cite{ref141,ref5}
that they cannot be quantized to get differential operators on any
line bundle other than $K_\cM^{1/2}$.   The fact that they can be
quantized to get {\it commuting} differential operators means that
Hitchin's classical integrable system on $\MH$ can be
``quantized'' to make a quantum integrable system on $\M$.
